\IfFileExists{stacks-project.cls}{%
\documentclass{stacks-project}
}{%
\documentclass{amsart}
}

% For dealing with references we use the comment environment
\usepackage{verbatim}
\newenvironment{reference}{\comment}{\endcomment}
%\newenvironment{reference}{}{}
\newenvironment{slogan}{\comment}{\endcomment}
\newenvironment{history}{\comment}{\endcomment}

% For commutative diagrams we use Xy-pic
\usepackage[all]{xy}

% We use 2cell for 2-commutative diagrams.
\xyoption{2cell}
\UseAllTwocells

% We use multicol for the list of chapters between chapters
\usepackage{multicol}

% This is generally recommended for better output
\usepackage{lmodern}
\usepackage[T1]{fontenc}

% For cross-file-references
\usepackage{xr-hyper}

% Package for hypertext links:
\usepackage{hyperref}

% For any local file, say "hello.tex" you want to link to please
% use \externaldocument[hello-]{hello}
\externaldocument[introduction-]{introduction}
\externaldocument[conventions-]{conventions}
\externaldocument[sets-]{sets}
\externaldocument[categories-]{categories}
\externaldocument[topology-]{topology}
\externaldocument[sheaves-]{sheaves}
\externaldocument[sites-]{sites}
\externaldocument[stacks-]{stacks}
\externaldocument[fields-]{fields}
\externaldocument[algebra-]{algebra}
\externaldocument[brauer-]{brauer}
\externaldocument[homology-]{homology}
\externaldocument[derived-]{derived}
\externaldocument[simplicial-]{simplicial}
\externaldocument[more-algebra-]{more-algebra}
\externaldocument[smoothing-]{smoothing}
\externaldocument[modules-]{modules}
\externaldocument[sites-modules-]{sites-modules}
\externaldocument[injectives-]{injectives}
\externaldocument[cohomology-]{cohomology}
\externaldocument[sites-cohomology-]{sites-cohomology}
\externaldocument[dga-]{dga}
\externaldocument[dpa-]{dpa}
\externaldocument[sdga-]{sdga}
\externaldocument[hypercovering-]{hypercovering}
\externaldocument[schemes-]{schemes}
\externaldocument[constructions-]{constructions}
\externaldocument[properties-]{properties}
\externaldocument[morphisms-]{morphisms}
\externaldocument[coherent-]{coherent}
\externaldocument[divisors-]{divisors}
\externaldocument[limits-]{limits}
\externaldocument[varieties-]{varieties}
\externaldocument[topologies-]{topologies}
\externaldocument[descent-]{descent}
\externaldocument[perfect-]{perfect}
\externaldocument[more-morphisms-]{more-morphisms}
\externaldocument[flat-]{flat}
\externaldocument[groupoids-]{groupoids}
\externaldocument[more-groupoids-]{more-groupoids}
\externaldocument[etale-]{etale}
\externaldocument[chow-]{chow}
\externaldocument[intersection-]{intersection}
\externaldocument[pic-]{pic}
\externaldocument[weil-]{weil}
\externaldocument[adequate-]{adequate}
\externaldocument[dualizing-]{dualizing}
\externaldocument[duality-]{duality}
\externaldocument[discriminant-]{discriminant}
\externaldocument[derham-]{derham}
\externaldocument[local-cohomology-]{local-cohomology}
\externaldocument[algebraization-]{algebraization}
\externaldocument[curves-]{curves}
\externaldocument[resolve-]{resolve}
\externaldocument[models-]{models}
\externaldocument[functors-]{functors}
\externaldocument[equiv-]{equiv}
\externaldocument[pione-]{pione}
\externaldocument[etale-cohomology-]{etale-cohomology}
\externaldocument[proetale-]{proetale}
\externaldocument[relative-cycles-]{relative-cycles}
\externaldocument[more-etale-]{more-etale}
\externaldocument[trace-]{trace}
\externaldocument[crystalline-]{crystalline}
\externaldocument[spaces-]{spaces}
\externaldocument[spaces-properties-]{spaces-properties}
\externaldocument[spaces-morphisms-]{spaces-morphisms}
\externaldocument[decent-spaces-]{decent-spaces}
\externaldocument[spaces-cohomology-]{spaces-cohomology}
\externaldocument[spaces-limits-]{spaces-limits}
\externaldocument[spaces-divisors-]{spaces-divisors}
\externaldocument[spaces-over-fields-]{spaces-over-fields}
\externaldocument[spaces-topologies-]{spaces-topologies}
\externaldocument[spaces-descent-]{spaces-descent}
\externaldocument[spaces-perfect-]{spaces-perfect}
\externaldocument[spaces-more-morphisms-]{spaces-more-morphisms}
\externaldocument[spaces-flat-]{spaces-flat}
\externaldocument[spaces-groupoids-]{spaces-groupoids}
\externaldocument[spaces-more-groupoids-]{spaces-more-groupoids}
\externaldocument[bootstrap-]{bootstrap}
\externaldocument[spaces-pushouts-]{spaces-pushouts}
\externaldocument[spaces-chow-]{spaces-chow}
\externaldocument[groupoids-quotients-]{groupoids-quotients}
\externaldocument[spaces-more-cohomology-]{spaces-more-cohomology}
\externaldocument[spaces-simplicial-]{spaces-simplicial}
\externaldocument[spaces-duality-]{spaces-duality}
\externaldocument[formal-spaces-]{formal-spaces}
\externaldocument[restricted-]{restricted}
\externaldocument[spaces-resolve-]{spaces-resolve}
\externaldocument[formal-defos-]{formal-defos}
\externaldocument[defos-]{defos}
\externaldocument[cotangent-]{cotangent}
\externaldocument[examples-defos-]{examples-defos}
\externaldocument[algebraic-]{algebraic}
\externaldocument[examples-stacks-]{examples-stacks}
\externaldocument[stacks-sheaves-]{stacks-sheaves}
\externaldocument[criteria-]{criteria}
\externaldocument[artin-]{artin}
\externaldocument[quot-]{quot}
\externaldocument[stacks-properties-]{stacks-properties}
\externaldocument[stacks-morphisms-]{stacks-morphisms}
\externaldocument[stacks-limits-]{stacks-limits}
\externaldocument[stacks-cohomology-]{stacks-cohomology}
\externaldocument[stacks-perfect-]{stacks-perfect}
\externaldocument[stacks-introduction-]{stacks-introduction}
\externaldocument[stacks-more-morphisms-]{stacks-more-morphisms}
\externaldocument[stacks-geometry-]{stacks-geometry}
\externaldocument[moduli-]{moduli}
\externaldocument[moduli-curves-]{moduli-curves}
\externaldocument[examples-]{examples}
\externaldocument[exercises-]{exercises}
\externaldocument[guide-]{guide}
\externaldocument[desirables-]{desirables}
\externaldocument[coding-]{coding}
\externaldocument[obsolete-]{obsolete}
\externaldocument[fdl-]{fdl}
\externaldocument[index-]{index}

% Theorem environments.
%
\theoremstyle{plain}
\newtheorem{theorem}[subsection]{Theorem}
\newtheorem{proposition}[subsection]{Proposition}
\newtheorem{lemma}[subsection]{Lemma}

\theoremstyle{definition}
\newtheorem{definition}[subsection]{Definition}
\newtheorem{example}[subsection]{Example}
\newtheorem{exercise}[subsection]{Exercise}
\newtheorem{situation}[subsection]{Situation}

\theoremstyle{remark}
\newtheorem{remark}[subsection]{Remark}
\newtheorem{remarks}[subsection]{Remarks}

\numberwithin{equation}{subsection}

% Macros
%
\def\lim{\mathop{\mathrm{lim}}\nolimits}
\def\colim{\mathop{\mathrm{colim}}\nolimits}
\def\Spec{\mathop{\mathrm{Spec}}}
\def\Hom{\mathop{\mathrm{Hom}}\nolimits}
\def\Ext{\mathop{\mathrm{Ext}}\nolimits}
\def\SheafHom{\mathop{\mathcal{H}\!\mathit{om}}\nolimits}
\def\SheafExt{\mathop{\mathcal{E}\!\mathit{xt}}\nolimits}
\def\Sch{\mathit{Sch}}
\def\Mor{\mathop{\mathrm{Mor}}\nolimits}
\def\Ob{\mathop{\mathrm{Ob}}\nolimits}
\def\Sh{\mathop{\mathit{Sh}}\nolimits}
\def\NL{\mathop{N\!L}\nolimits}
\def\CH{\mathop{\mathrm{CH}}\nolimits}
\def\proetale{{pro\text{-}\acute{e}tale}}
\def\etale{{\acute{e}tale}}
\def\QCoh{\mathit{QCoh}}
\def\Ker{\mathop{\mathrm{Ker}}}
\def\Im{\mathop{\mathrm{Im}}}
\def\Coker{\mathop{\mathrm{Coker}}}
\def\Coim{\mathop{\mathrm{Coim}}}

% Boxtimes
%
\DeclareMathSymbol{\boxtimes}{\mathbin}{AMSa}{"02}

%
% Macros for moduli stacks/spaces
%
\def\QCohstack{\mathcal{QC}\!\mathit{oh}}
\def\Cohstack{\mathcal{C}\!\mathit{oh}}
\def\Spacesstack{\mathcal{S}\!\mathit{paces}}
\def\Quotfunctor{\mathrm{Quot}}
\def\Hilbfunctor{\mathrm{Hilb}}
\def\Curvesstack{\mathcal{C}\!\mathit{urves}}
\def\Polarizedstack{\mathcal{P}\!\mathit{olarized}}
\def\Complexesstack{\mathcal{C}\!\mathit{omplexes}}
% \Pic is the operator that assigns to X its picard group, usage \Pic(X)
% \Picardstack_{X/B} denotes the Picard stack of X over B
% \Picardfunctor_{X/B} denotes the Picard functor of X over B
\def\Pic{\mathop{\mathrm{Pic}}\nolimits}
\def\Picardstack{\mathcal{P}\!\mathit{ic}}
\def\Picardfunctor{\mathrm{Pic}}
\def\Deformationcategory{\mathcal{D}\!\mathit{ef}}
\usepackage{cleveref}
\crefname{originalreference}{original reference}{original references}
\Crefname{originalreference}{Original reference}{Original references}
\makeatletter
\newlabel{bootstrap-lemma-descend-algebraic-space}{{0ADV}{source}{}{sourcepointer-1}{}}
\newlabel{bootstrap-lemma-descend-algebraic-space@cref}{{[originalreference][0][]0ADV}{[1][0][]source}{}{}{}}
\makeatother
\crefname{originalequation}{original equation}{original equations}
\Crefname{originalequation}{Original equation}{Original equations}
\makeatletter
\newlabel{equation-completion-functor}{{0AR5}{source}{}{sourcepointer-2}{}}
\newlabel{equation-completion-functor@cref}{{[originalequation][0][]0AR5}{[1][0][]source}{}{}{}}
\makeatother
\crefname{originallemma}{original lemma}{original lemmas}
\Crefname{originallemma}{Original lemma}{Original lemmas}
\makeatletter
\newlabel{lemma-completion-functor-fully-faithful}{{0GDR}{source}{}{sourcepointer-3}{}}
\newlabel{lemma-completion-functor-fully-faithful@cref}{{[originallemma][0][]0GDR}{[1][0][]source}{}{}{}}
\makeatother
\crefname{originallemma}{original lemma}{original lemmas}
\Crefname{originallemma}{Original lemma}{Original lemmas}
\makeatletter
\newlabel{lemma-dilatations-affine}{{0ARA}{source}{}{sourcepointer-4}{}}
\newlabel{lemma-dilatations-affine@cref}{{[originallemma][0][]0ARA}{[1][0][]source}{}{}{}}
\makeatother
\crefname{originallemma}{original lemma}{original lemmas}
\Crefname{originallemma}{Original lemma}{Original lemmas}
\makeatletter
\newlabel{lemma-functoriality-completion-functor}{{0GDQ}{source}{}{sourcepointer-5}{}}
\newlabel{lemma-functoriality-completion-functor@cref}{{[originallemma][0][]0GDQ}{[1][0][]source}{}{}{}}
\makeatother
\crefname{originalreference}{original reference}{original references}
\Crefname{originalreference}{Original reference}{Original references}
\makeatletter
\newlabel{spaces-descent-definition-descent-datum-for-family-of-morphisms}{{0ADI}{source}{}{sourcepointer-6}{}}
\newlabel{spaces-descent-definition-descent-datum-for-family-of-morphisms@cref}{{[originalreference][0][]0ADI}{[1][0][]source}{}{}{}}
\makeatother
\title{The completion functor is an equivalence on rig-etale morphisms}
\author{The Stacks Project Authors}
\date{stacks-project:restricted.tex:local-snapshot-20261012; GNU FDL 1.2 or later}
\begin{document}
\maketitle
\noindent\textit{Editorial note.} Counted span is the author's own proof environment, source lines 7042-7097, following the statement at 7014-7040.\par\medskip
\noindent\textit{Editorial note.} The final 'reader can check ... some details omitted' step is the sheaf-level gluing of already-established descent data, flagged in the notes.\par\medskip
\noindent\textit{Editorial note.} Steps the author leaves to the reader are flagged here for the reviewer: source line(s) 7096 carry an "omitted"/"follows easily"/"immediate from the definitions" phrase. Those sentences are retained verbatim; the displayed argument around them stays inside the counted span.\par\medskip
\noindent\textit{Allowed standard prerequisites (editorial).}\begin{itemize}
\item Fully faithfulness of the completion functor. lemma-completion-functor-fully-faithful (carded L26155, cited at \textasciitilde{}7043). Twice: for the a\_ij and their cocycle condition.
\item Affine dilatations. lemma-dilatations-affine (carded L26272, cited at \textasciitilde{}7055). Supplies X'\_i -> X\_i with (X'\_i)\_\{/T\} = W\_i.
\item Effectiveness of descent data. Descent on Spaces, Definition spaces-descent-definition-descent-datum-for-family-of-morphisms and Bootstrap, Lemma bootstrap-lemma-descend-algebraic-space (cited at \textasciitilde{}7072-7076). Descends the X'\_i to X'.
\end{itemize}\medskip
\ifdefined\hypertarget\hypertarget{sourcepointer-1}{}\fi
\noindent\textit{Source reference.} Original reference 0ADV: focused-seeds/geometry/S9-restricted-dilatations-general/ref-tags.txt, source lines 5--5. Navigation pointer only: the counted proof cites this statement by \textbackslash{}ref as a standing input; the deduction the counted proof has to make is carried out inside the counted span, and no part of that deduction is delegated to this pointer.\par\medskip
\ifdefined\hypertarget\hypertarget{sourcepointer-2}{}\fi
\noindent\textit{Source reference.} Original equation 0AR5: sources/stacks-restricted/pinned-original-source.tex, source lines 6766--6766. Navigation pointer only: the counted proof cites this statement by \textbackslash{}ref as a standing input; the deduction the counted proof has to make is carried out inside the counted span, and no part of that deduction is delegated to this pointer.\par\medskip
\ifdefined\hypertarget\hypertarget{sourcepointer-3}{}\fi
\noindent\textit{Source reference.} Original lemma 0GDR: sources/stacks-restricted/pinned-original-source.tex, source lines 6832--6832. Navigation pointer only: the counted proof cites this statement by \textbackslash{}ref as a standing input; the deduction the counted proof has to make is carried out inside the counted span, and no part of that deduction is delegated to this pointer.\par\medskip
\ifdefined\hypertarget\hypertarget{sourcepointer-4}{}\fi
\noindent\textit{Source reference.} Original lemma 0ARA: sources/stacks-restricted/pinned-original-source.tex, source lines 6946--6946. Navigation pointer only: the counted proof cites this statement by \textbackslash{}ref as a standing input; the deduction the counted proof has to make is carried out inside the counted span, and no part of that deduction is delegated to this pointer.\par\medskip
\ifdefined\hypertarget\hypertarget{sourcepointer-5}{}\fi
\noindent\textit{Source reference.} Original lemma 0GDQ: sources/stacks-restricted/pinned-original-source.tex, source lines 6790--6790. Navigation pointer only: the counted proof cites this statement by \textbackslash{}ref as a standing input; the deduction the counted proof has to make is carried out inside the counted span, and no part of that deduction is delegated to this pointer.\par\medskip
\ifdefined\hypertarget\hypertarget{sourcepointer-6}{}\fi
\noindent\textit{Source reference.} Original reference 0ADI: focused-seeds/geometry/S9-restricted-dilatations-general/ref-tags.txt, source lines 6--6. Navigation pointer only: the counted proof cites this statement by \textbackslash{}ref as a standing input; the deduction the counted proof has to make is carried out inside the counted span, and no part of that deduction is delegated to this pointer.\par\medskip
\begin{theorem}
\label{theorem-dilatations-general}
Let $S$ be a scheme. Let $X$ be a locally Noetherian algebraic space over $S$.
Let $T \subset |X|$ be a closed subset. Let $U \subset X$ be the open subspace
with $|U| = |X| \setminus T$. The completion functor
(\ref{equation-completion-functor})
$$
\left\{
\begin{matrix}
\text{morphisms of algebraic spaces}\\
f : X' \to X\text{ which are locally}\\
\text{of finite type and such that}\\
f^{-1}U \to U\text{ is an isomorphism}
\end{matrix}
\right\}
\longrightarrow
\left\{
\begin{matrix}
\text{morphisms }g : W \to X_{/T}\\
\text{of formal algebraic spaces}\\
\text{with }W\text{ locally Noetherian}\\
\text{and }g\text{ rig-\'etale}
\end{matrix}
\right\}
$$
sending $f : X' \to X$ to $f_{/T} : X'_{/T'} \to X_{/T}$ is an equivalence.
\end{theorem}

\begin{proof}
The functor is fully faithful by
Lemma \ref{lemma-completion-functor-fully-faithful}.
Let $g : W \to X_{/T}$ be a morphism of formal algebraic spaces
with $W$ locally Noetherian and $g$ rig-\'etale.
We will prove $W$ is in the essential image to finish the proof.

\medskip\noindent
Choose an \'etale covering $\{X_i \to X\}$ with $X_i$ affine for all $i$.
Denote $U_i \subset X_i$ the inverse image of $U$ and denote
$T_i \subset X_i$ the inverse image of $T$.
Recall that $(X_i)_{/T_i} = (X_i)_{/T} = (X_i \times_X X)_{/T}$ and
$W_i = X_i \times_X W = (X_i)_{/T} \times_{X_{/T}} W$, see
Lemma \ref{lemma-functoriality-completion-functor}.
Observe that we obtain isomorphisms
$$
\alpha_{ij} :
W_i \times_{X_{/T}} (X_j)_{/T}
\longrightarrow 
(X_i)_{/T} \times_{X_{/T}} W_j
$$
satisfying a suitable cocycle condition.
By Lemma \ref{lemma-dilatations-affine} applied to
$X_i, T_i, U_i, W_i \to (X_i)_{/T}$
there exists a morphism $X'_i \to X_i$ of algebraic spaces
which is locally of finite type and an isomorphism over $U_i$
and an isomorphism $\beta_i : (X'_i)_{/T} \cong W_i$ over $(X_i)_{/T}$.
By fully faithfullness we find an isomorphism
$$
a_{ij} : X'_i \times_X X_j \longrightarrow X_i \times_X X'_j
$$
over $X_i \times_X X_j$ such that
$\alpha_{ij} = \beta_j|_{X_i \times_X X_j}
\circ (a_{ij})_{/T} \circ \beta_i^{-1}|_{X_i \times_X X_j}$.
By fully faithfulness again (this time over
$X_i \times_X X_j \times_X X_k$)
we see that these morphisms $a_{ij}$ satisfy the same
cocycle condition as satisfied by the $\alpha_{ij}$.
In other words, we obtain a descent datum
(as in Descent on Spaces, Definition
\ref{spaces-descent-definition-descent-datum-for-family-of-morphisms})
$(X'_i, a_{ij})$ relative to the family $\{X_i \to X\}$.
By Bootstrap, Lemma \ref{bootstrap-lemma-descend-algebraic-space},
this descent datum is effective. Thus we find a morphism
$f : X' \to X$ of algebraic spaces and isomorphisms
$h_i : X' \times_X X_i \to X'_i$ over $X_i$ such that
$a_{ij} = h_j|_{X_i \times_X X_j} \circ h_i^{-1}|_{X_i \times_X X_j}$.
The reader can check that the ensuing isomorphisms
$$
(X' \times_X X_i)_{/T}
\xrightarrow{\beta_i \circ (h_i)_{/T}}
W_i
$$
over $X_i$ glue to an isomorphism $X'_{/T} \to W$
over $X_{/T}$; some details omitted.
\end{proof}

\end{document}
